\documentclass[../main.tex]{subfiles}
\begin{document}
%==========================================TIÊU ĐỀ TẬP========================================
\begin{table}[h]
    \begin{adjustwidth}{-1cm}{}
        \begin{tabular}{>{\centering\arraybackslash}p{6cm}|>{\raggedright\arraybackslash}p{14cm}}
            \multirow{5}{6cm}
            % Cột bên trái: ÉPISODE và số tập
            {\\[-40pt]
                \flaregothic\fontsize{20pt}{20pt}\selectfont \centering
                \color{eptype}ÉPISODE\color{black}\\
                \burbank\fontsize{72pt}{72pt}\selectfont
                \color{epnum}6\color{black}\\[6pt]
                \cafeteria\fontsize{20pt}{20pt}\selectfont\color{epnum}(1-2)\color{black}
                \\[18pt]
                % \flaregothic\fontsize{14pt}{14pt}\selectfont
                % \color{epattr}BIENVENUE, \uline{\color{Banpire} Mel} !
                \flaregothic\fontsize{18pt}{18pt}\selectfont
                \color{epattr}ÉPISODE PERDU
                \color{black}
            }
             &
            % Dòng thứ nhất: tên tập tiếng Nhật
            {\fotskip\fontsize{18pt}{18pt}\selectfont \ruby{JK}{女子高生}でもクマに\ruby{勝}{か}てるライフハック
            }                                                                                       \\[6pt]
            % Dòng thứ hai: tên tập tiếng Anh
             & {\cafeteria\fontsize{18pt}{18pt}\selectfont TEEN Beating BEAR Life Hack!!}                                                                                          \\[4pt]
            % Dòng thứ ba: tên tập tiếng Việt
             & {\vnmsans\fontsize{18pt}{18pt}\selectfont Chụp ảnh với gấu}                                                                                                         \\[8pt]
            % Dòng thứ tư: Nhân vật xuất hiện
             & {\caviar{\fontsize{14pt}{14pt}\selectfont Nhân vật: \emp{kuon1} \emp{roboco} \emp{matsuri} \emp{haato} \emp{aqua} \emp{shion} \emp{ayame} \emp{choco} \emp{subaru}}} \\[8pt]
            % Dòng cuối cùng: Thời gian diễn ra của tập (thường sẽ là ngày phát hành)
             & {\continuum\fontsize{16pt}{16pt}\selectfont Ngày phát hành: 09/06/2019}                                                                                             \\[18pt]
        \end{tabular}
    \end{adjustwidth}
\end{table}

\color{black}

\txtbx[2]{
    {\color{vol} \uline{Tóm tắt:}} Để chuẩn bị cho buổi chụp ảnh quảng bá của công ty, Haato vô tình sử dụng một chiếc máy ảnh bị nguyền rủa của gia đình Ayame, thứ hút hồn bất kỳ ai có bức ảnh không đạt yêu cầu. Sau khi Matsuri suýt nữa mất mạng và được cứu thoát, cả nhóm nỗ lực tìm kiếm trang phục đẹp để làm hài lòng chiếc máy ảnh, nhưng lại tiếp tục đụng độ bộ đồ gấu bông bị nguyền rủa khác biến Roboco thành một chú gấu đáng yêu. Buổi chụp hình đầy hỗn loạn kết thúc khi chú gấu Roboco hồn lìa khỏi xác trước vẻ dễ thương khó cưỡng của Haato, để lại những tấm ảnh ngộ nghĩnh làm kỷ niệm cho Kuon và các bạn cùng lớp CS1.
}

Hôm nay, văn phòng \hll\ nhộn nhịp hơn hẳn ngày thường bởi một buổi chụp ảnh quảng bá toàn công ty. Đây là một trong những cơ hội hiếm hoi để các cô gái nhà Tam giác Xanh thỏa sức tạo dáng trước ống kính, chuẩn bị những bức hình lung linh cho các sự kiện sắp tới. Mục tiêu không chỉ dừng lại ở thị trường trong nước, mà còn hướng tới việc vươn tầm quốc tế, đúng như giấc mơ cháy bỏng về một thế hệ idol toàn cầu của ngài CEO Tanigo.

Tuy vậy, Kuon thì vốn chẳng mấy mặn mà với việc chụp ảnh. Cậu chọn cho mình một góc khuất yên tĩnh trên chiếc ghế sô-pha quen thuộc, lặng lẽ quan sát các cô gái đang tất bật làm việc. Cậu không muốn làm gián đoạn guồng quay bận rộn của họ, và càng không muốn can thiệp sâu vào một lĩnh vực vốn chẳng phải là sở trường của mình.

\dots Dẫu rằng trước đây cậu cũng có chút tiếng tăm nhờ những bức ảnh chụp gia đình ấm áp, nhưng chụp ảnh chuyên nghiệp chưa bao giờ là niềm đam mê thực sự của cậu. ``\textit{Để xem mấy đứa nhóc này làm ăn ra sao,}'' cậu thầm nghĩ, đưa cốc nước đang uống dở lên nhấp một ngụm nhỏ đầy thư thái.

\img{1-2a}

Haato thích thú nâng chiếc máy ảnh hiệu Holon\footnote{Nhái của Nikon.} - thứ mà cô nàng vô tình lượm được trên bàn của ai đó bỏ quên - ra hiệu cho mọi người tập trung. Cô lùi lại vài bước để căn góc, giọng lảnh lót: ``Mọi người ơi\~{} Xếp hàng nào, xếp hàng nào! Chuẩn bị chụp ảnh quảng bá thôi!''

Matsuri lập tức tiến lên trước mặt nàng tóc vàng, đưa hai tay ra trước ngực, khẽ nắm lại đầy nhí nhảnh cùng một nụ cười rạng rỡ nở trên môi. Bên cạnh cô, Roboco đang ngồi thong thả ở phía bên trái, trong khi Kuon vẫn duy trì khoảng cách ở góc phòng xa xa.

``Tươi lên nào!'' Haato nói, và một tiếng \textbf{*TÁCH!*} vang lên từ chiếc máy ảnh mà cô cầm trên tay, cùng với một ánh chớp lóe lên khắp phòng. Cô nàng hạ máy xuống, đang định bảo đối phương đổi dáng thì bất chợt cô cùng với những người khác nhìn thấy một hiện tượng lạ.

\img{1-2b}

``Tạm biệt mọi người\dots\ Cảm ơn vì tất cả\dots'' toàn thân cô nàng Cổ động cứng đờ lại, chuyển sang màu xám xịt. Con ngươi trên mắt cô trợn ngược lên trên, cùng với phần hồn (?) của cô bắt đầu lìa khỏi xác.

``C-Cái quái gì thế này!?'' Kuon thốt lên kinh ngạc, suýt nữa thì đánh rơi cốc nước đang uống dở.

Cả Haato và Roboco đều đồng thanh hét lên đầy kinh hãi: ``\textbf{Matsuri-chan!?}''

``Này, chuyện gì đang xảy ra vậy!?'' cậu lập tức bật dậy khỏi ghế sô-pha, chỉ tay về phía linh hồn đang bay lơ lửng của cô nàng. ``Em định bay đi đầu thai cùng chú bé liên lạc\footnote{Ý nói đến nhân vật Lượm trong bài thơ cùng tên của Tố Hữu. Trong bài thơ, cậu bé Lượm đã hy sinh trong khi làm nhiệm vụ liên lạc khi tuổi đời còn rất nhỏ.} luôn đấy à!?''

``Matsuri à!? Matsuri-chan!!'' Haato hốt hoảng khóc ròng, giơ hai tay cố túm lấy cái bóng xám xịt của người bạn.

``\textbf{Đứng lại đó!}'' Một tiếng hét dõng dạc bất chợt vang lên từ phía cửa ra vào. Ayame từ đâu trượt thẳng vào văn phòng, hai tay dang rộng chuẩn bị tư thế nghênh chiến như một dũng sĩ thực thụ. Sự xuất hiện bất ngờ của nàng Quỷ sừng khiến cả Haato và Roboco chỉ biết tròn mắt đồng thanh gọi tên cô, trong khi cậu Tổng Kuon vẫn còn đang ngơ ngác chưa hiểu mô tê gì.

Không chậm trễ một giây, cô lao vút tới như một cơn gió. Cô vươn tay chộp lấy phần linh hồn xám xịt đang chuẩn bị tan biến của Matsuri, dồn lực đập thẳng nó trở lại đầu của cô nàng. Một luồng sáng nhẹ lóe lên, kéo cái bóng xám xịt nhập lại vào thể xác, đưa Matsuri quay trở về trạng thái bình thường trong gang tấc.

Nàng Cổ động lảo đảo đứng vững, lồng ngực phập phồng hú vía sau trải nghiệm kinh hoàng vừa trải qua. Cảm giác lúc này của cô chẳng khác nào vừa từ cõi chết trở về. Đứng bên cạnh cô, nàng tóc trắng Ayame thở phào nhẹ nhõm, đưa tay quệt những giọt mồ hôi lấm tấm trên trán.

Matsuri khuỵu gối xuống sàn, ôm lấy ngực và thở hồng hộc đầy nhọc nhằn: ``Em\dots\ em vừa nhìn thấy dòng sông Styx sao\dots\footnote{Theo thần thoại Hy Lạp, đây là một con sông tạo nên ranh giới giữa trần gian và âm phủ. Ý của Matsuri là cô vừa trở về từ cõi chết.}''

``Em suýt nữa là lên bàn thờ ngắm gà khỏa thân rồi đó,'' Kuon thở phào một hơi rồi quay sang nhìn nàng Nữ sinh trung học hỏi tội. ``Nhưng tại sao tự dưng em lại bị hồn bay phách lạc như thế? Haato-chan, em vừa làm gì à?''

``Đâu có!'' Haato giơ hai tay lên phân bua, khuôn mặt méo xệch vì oan uổng. ``Em thề là em chỉ mới bấm nút chụp hình bình thường thôi mà!''

Ayame liền vội vã giải thích về thứ quái dị mà nàng tóc vàng đang cầm trên tay bằng giọng run rẩy: ``Đó là chiếc máy ảnh bị nguyền rủa được truyền thừa qua nhiều thế hệ của gia tộc em đấy! Nếu bức hình chụp ra không đạt yêu cầu thẩm mỹ của chiếc máy, linh hồn của người được chụp sẽ bị hút đi ngay lập tức!''

\img{1-2c}

Haato và Roboco nhìn chiếc máy ảnh cơ cũ kỹ trên bàn với ánh mắt đầy lo âu. Vỏ ngoài của nó bỗng chốc tỏa ra thứ ánh sáng xanh tím kỳ ảo nhưng lạnh lẽo, ma mị đến rợn người. Kuon cau mày bước tới thắc mắc: ``Vậy ra thứ đồ chơi nguy hiểm này là của em à? Sao em lại vứt cái của nợ này ở văn phòng chứ!?''

``Xin lỗi anh, Kuon-senpai! Tại em vội quá cho nên quên cất đi\dots'' Ayame lí nhí cúi đầu, khuôn mặt hoang mang tột độ trước thứ sức mạnh ma quái đang rục rịch hoạt động. ``Có lẽ bức ảnh của Matsuri-senpai lúc nãy chụp không đạt chuẩn thẩm mỹ của nó, nên mới dẫn đến hiện tượng hút hồn đáng sợ kia đấy ạ!''

Nghe đến đây, cả phòng đều rùng mình ớn lạnh trước sự tồn tại của thứ đồ vật quỷ dị này. Thế nhưng, ngoại trừ nạn nhân suýt mất mạng là Matsuri, cô nàng dường như chẳng hề để tâm đến nỗi sợ hãi đó. Đôi mắt nàng Cổ động bỗng sáng lên, cô đập hai tay vào nhau reo lên: ``Em có ý này! Nếu máy ảnh đòi hỏi thẩm mỹ, sao chúng ta không mặc những bộ đồ thật đẹp để chụp ảnh nhỉ?''

Cậu Tổng và nàng Quỷ sừng đồng loạt nhìn cô với ánh mắt đầy hoài nghi, trong đầu cùng chung một suy nghĩ: ``\textit{Em ấy bị làm sao thế này? / Đúng là điếc không sợ súng mà.}'' Họ thầm hy vọng Matsuri chỉ là người duy nhất có cái ý tưởng phiêu lưu điên rồ kia--

``Tuyệt vời! Chúng ta cùng đi tìm những bộ trang phục lộng lẫy nhất nào!'' Roboco vui vẻ hưởng ứng, lập tức hùa theo ý tưởng đó. Haato đứng bên cạnh đổ mồ hôi hột, ái ngại nhìn hai người họ: ``Nhưng mà\dots\ đi kiếm một cái máy ảnh bình thường không phải là nhanh và an toàn hơn sao?''

Cậu con trai gật gù đồng tình: ``Cái đó thì anh đồng ý cả hai tay hai chân luôn. Mau tìm cho anh cái máy ảnh nào bình thường một chút đi--''

Bất chợt, một đoạn quảng cáo giới thiệu các thành viên Thế hệ thứ Hai gồm Choco, Shion, Ayame, Subaru và Aqua\footnote{Ứng với năm chị em gái nhà Nakano trong \textit{Nhà có năm nàng dâu - \ruby{五}{ご}\ruby{等}{とう}\ruby{分}{ぶん}の\ruby{花}{はな}\ruby{嫁}{よめ}}, từ lớn nhất đến bé nhất: Ichika - Nino - Miku - Yotsuba và Itsuki.}. Cả năm cô gái đứng sát vai nhau, ống kính máy quay lướt nhẹ một vòng khoe trọn tạo hình độc đáo của từng người trong thế hệ, kèm theo đó là giọng thuyết minh ngọt ngào nhưng đầy kịch tính của Aqua giới thiệu về bản đặc biệt thương hiệu.

\gif{1-2d}{1}{161}

``Khoan, khoan, khoan!'' Kuon đứng bật dậy, hai tay gạt ngang mặt như muốn xua tan cái quảng cáo vô lý vừa rồi. ``Từ khi nào mà mấy đứa tự tiện đi làm quảng cáo vậy hả!? Với lại, tại sao lại nhái cái bộ anime đó!?''

Bất chấp sự can ngăn của cậu Tổng, Aqua trong đoạn băng quảng cáo vẫn tươi cười rạng rỡ, tay cầm hộp đĩa giơ lên trước ống kính: ``\textit{Nhà có 5 thành viên thế hệ 2} giờ đã có bản Blu-ray và DVD! Mọi người hãy nhanh tay sở hữu chúng nhé!''

Cậu Tổng bất lực gục đầu xuống bàn, thở dài ngao ngán: ``\dots\ Anh nghiêm túc đấy. Mấy đứa không sợ bị bên bản quyền kiện cho sập tiệm à?''

\ct

Sau khoảng một giờ đồng hồ lục tung mọi ngóc ngách của văn phòng để tìm kiếm những bộ cánh lộng lẫy nhất, cả nhóm mệt nhoài tụ tập lại ở giữa phòng để tổng kết thành quả.

Và khi nhắc đến ``cả nhóm'', dĩ nhiên bao gồm cả cậu Tổng Kuon. Dù bản thân chẳng hề muốn dính líu vào cái trò hóa trang đầy rủi ro này, cậu vẫn đành ngậm ngùi tham gia chỉ vì ``các cô nàng cứ nài nỉ bắt ép bằng được.''

Đáng tiếc thay, tất cả đống trang phục thu thập được rốt cuộc chỉ có độc nhất một bộ đồ hóa trang con gấu nâu, xộc xệch mà chính tay cậu Tổng vừa mang từ phòng làm việc của ngài CEO Tanigo sang.

``Chúng ta rõ ràng là một công ty quản lý idol tầm cỡ,'' Haato nhăn mặt bưng bộ đồ gấu bông lên, giọng điệu đầy bất mãn, ``Thế mà trong kho chỉ có độc nhất một bộ đồ Ngài M'goosefish-Liver này thôi sao!?''

Kuon nhún vai, nhấp một ngụm cốc nước rồi thản nhiên đáp: ``Chắc tại phía đối tác thấy bọn em giống một đám hề xiếc hơn là idol thực thụ, nên họ mới không thèm tốn công may thêm đồ đấy.''

``\textbf{Cái gì cơ!?}'' nàng Nữ sinh trung học đỏ mặt tía tai vì cáu, giậm chân bình bịch xuống sàn.

``Họ nhận xét thế chứ anh có bảo đâu,'' Kuon giơ hai tay lên cười trừ, khuôn mặt tỉnh bơ đầy trêu chọc. ``Với lại, ngẫm kỹ thì họ có đâu có sai.''

Cậu con trai nhìn chiếc máy ảnh của Ayame đặt ở trên bàn, tỏ vẻ nửa tin nửa ngờ về nó. ``Anh không nghĩ cái máy này thật sự nguyền rủa đến thế đâu,'' cậu nói, cầm chiếc máy ảnh lên và chụp Matsuri đang nói chuyện với Ayame.

Thế nhưng, cậu vừa hạ máy xuống, hồn của cô nàng Cổ động lại bắt đầu bay đi mất khiến cho nàng Quỷ sừng hốt hoảng và kéo nó trở lại. Ojou-san nhìn cậu, tỏ vẻ không hài lòng với sự táy máy của cậu, còn bản thân cậu chỉ có thể thốt lên bằng tiếng mẹ đẻ trong sự kinh ngạc: ``\vie{Vãi cả linh hồn\dots}''

Trong khi đó, Roboco xung phong mặc bộ đồ gấu đó lên. Haato nhìn cô nàng Người máy với sự bất an.

``Cơ mà chị lại thực sự thấy ổn khi mặc cái bộ đồ quái dị này sao, chị Roboco\dots'' Haato nhìn cô nàng người máy đang hớn hở chui đầu vào bộ trang phục với ánh mắt đầy bất an.

Roboco-san hí hửng vẫy vẫy hai cái tay gấu bông bồng bềnh, giọng ngân dài điệu đà: ``Bởi vì nó đáng yêu cực kỳ luôn mà\~{}''

``Thôi được rồi\dots'' cô nàng Trung học thở dài thườn thượt, lẩm nhẩm cầu nguyện. ``Em chỉ hy vọng là từ giờ đến cuối buổi chụp sẽ không có chuyện gì kỳ lạ xảy ra nữa.''

Và đó chính là khoảnh khắc mà chính cô cảm thấy hối hận khi cô nói ra những lời đó. Ngay khi nàng Người máy đội chiếc đầu con gấu đó lên, lập tức cả bộ đồ tỏa ra thứ ánh sáng hắc ám, giống với chiếc máy ảnh ban nãy. Nàng Quỷ sừng chỉ tay về phía bộ đồ đó và lại hoảng hốt.

\img{1-2e}

``Chết rồi!'' Ayame hốt hoảng hét lên, mặt cắt không còn giọt máu khi nhìn thấy làn khói đen tím bao quanh cơ thể đối phương. ``Bộ trang phục M'goosefish-Liver đó cũng bị dính lời nguyền ma quỷ mất rồi!''

``Đừng có bảo với anh là cái thứ này cũng do tổ tiên em truyền lại đấy nhé!?'' Kuon đứng bật dậy, trán nổi đầy gân xanh bực bội.

Ayame cúi đầu, hai ngón tay trỏ chọc chọc vào nhau đầy hối lỗi: ``E-Em cũng không muốn thừa nhận đâu, nhưng\dots\ đúng là đồ gia bảo nhà em thật ạ.''

``Sao trong nhà em toàn chứa mấy cái thứ tà đạo nguy hiểm thế hả!?'' Kuon ôm đầu ngửa mặt lên trần nhà kêu trời, giọng điệu đầy bất lực: ``\textbf{\color{blue} Tại sao cơ chứ!?}''

``E-Em thật sự không nhớ là nó có tồn tại trong kho đồ mà!'' nàng Quỷ sừng xua hai tay giải thích trong sự hoảng loạn.

Bộ đồ gấu kỳ dị ấy bằng một cách thần bí nào đó đã biến Roboco thành một con gấu xám khổng lồ thực sự. Sinh vật to lớn ấy gầm rú lên những tiếng vang dội khắp phòng, khiến Ayame sợ đến mức mặt cắt không còn giọt máu, chỉ biết ôm đầu run lẩy bẩy. Matsuri thì đầu óc vẫn còn lâng lâng, chưa hoàn toàn hồi phục sau sự cố hồn bay phách lạc lúc nãy nên ngơ ngác nhìn quanh. Riêng chỉ có Haato và Kuon là vẫn duy trì được vẻ mặt bình tĩnh đến lạ thường trước cảnh tượng siêu thực này.

``Eeek! Đ-Đừng có ăn thịt em mà\dots'' Ayame co rúm người lại sau lưng Kuon, run lẩy bẩy cầu xin.

Matsuri thì mắt vẫn còn lờ đờ, đầu óc quay cuồng sau vụ suýt chết ban nãy, chỉ biết thì thào: ``Nếu đây thực sự là Roboco-senpai thì\dots''

``Thật không thể ngờ nổi luôn ấy,'' Haato lắc đầu chép miệng, đứng khoanh tay quan sát sinh vật khổng lồ trước mặt.

``Chuẩn cơm mẹ nấu luôn,'' Kuon chêm vào một câu đầy đồng tình.

Con gấu khổng lồ ấy lại há miệng gầm lên một lần nữa. Thế nhưng, thay vì một tiếng gầm rú dũng mãnh của chúa tể rừng xanh làm rúng động lòng người, từ trong họng của nó lại phát ra một âm thanh\dots\ vô cùng lạc quẻ và kỳ quặc.

``Meo\~{}!''

Bầu không khí trong văn phòng lập tức rơi vào trạng thái im lặng đến đóng băng. Tất cả mọi người đều đờ đẫn cả người, ngơ ngác nhìn chú gấu khổng lồ đang cố tỏ ra thân thiện trước mắt.

``C-Cái quái gì thế này?'' Kuon trố mắt ra nhìn, cốc nước trên tay suýt nữa đổ ụp xuống quần. ``Roboco-san? Em\dots\ em vẫn là em đấy chứ?''

``Ủa, đây có thực sự là tiếng kêu của loài gấu không vậy mọi người?'' Haato nghếch đầu thắc mắc.

Bất chợt, con gấu khổng lồ cúi xuống nhìn cô nàng tóc vàng, phát ra chất giọng robot quen thuộc của Roboco nhưng có phần trầm hơn: ``Thế em có giỏi thì gầm thử tiếng gấu cho chị xem nào.''

Nàng tóc vàng sửng sốt với yêu cầu của nó. Chưa nói đến việc nó (?) có thể nói như tiếng người. ``Ể!?'' cô thốt lên, còn Kuon thì thở dài: ``Anh mới là người cần yêu cầu em gầm như tiếng gấu như ban nãy đó\dots''

Haato chưa biết phải làm thế nào, đành phải tạo dáng và kêu lên một thứ tiếng mà như Kuon nói sau đó là ``nghe chả giống gấu một tí nào'.'

\img{1-2f}

``G-Gầm\dots\ Gào\dots\ Như thế này đã được chưa ạ?'' Haato đỏ mặt, hai tay nắm lại đưa lên má rồi khẽ kêu lên một tiếng đầy ngượng ngùng.

Kuon chứng kiến cảnh đó thì trán đổ môi hột, chỉ biết thầm nghĩ: ``\textit{Làm cái trò gì thế này\dots}''

Nhưng Matsuri thì hai mắt sáng rực lên, nhảy cẫng lên vỗ tay phấn khích: ``Nghe cũng ra dáng gấu đấy chứ! Đáng yêu quá đi mất thôi\~{}''

``\dots\ Cũng đúng nhỉ,'' Ayame gật gù phụ họa.

\img{1-2g}

Nhưng điều mà không ai ngờ tới chính là linh hồn của Roboco đột ngột thoát xác, bay lơ lửng ra khỏi bộ đồ gấu. Một tiếng kẻng ma quái vang lên vang dội. Mọi người trố mắt nhìn linh hồn xám xịt đang từ biệt cõi trần, trong khi lời trăng trối cuối cùng của cô nàng trước khi chính thức ``lìa đời'' chỉ là: ``Tôi chết vì sự dễ thương này mất thôi\dots''

Ngay sau đó, bộ đồ con gấu bị nguyền rủa bỗng tự bốc cháy rồi hóa thành tro bụi màu xám, không để lại bất kỳ dấu vết nào. Những người còn lại trong phòng há hốc mồm kinh ngạc, hoàn toàn cạn lời trước chuỗi sự kiện kỳ quặc vừa diễn ra trước mắt.

``Cái gì\dots\ vừa xảy ra thế này?'' tất cả bọn họ đồng thanh thốt lên trong sự ngơ ngác tột độ.

Quả thật, đây là một buổi chụp ảnh quảng bá hỗn loạn, kỳ bí và tràn ngập những tai nạn dở khóc dở cười nhất từ trước đến nay!

%==========================================TRUYỆN PHỤ=========================================
\omk

\begin{center}
    {\textit{*Các nhân vật nói bằng tiếng Etsunan.}}
\end{center}

\pov{Hikawa Kuon}

Vài ngày sau buổi chụp hình bão táp đầy rẫy những hiện tượng siêu thực ấy, Ayame đã phải lập tức thu dọn toàn bộ đống đồ gia bảo ma quái của mình ra khỏi văn phòng. Nếu không, em ấy chắc chắn sẽ phải đối mặt với cơn thịnh nộ của A-san, và tệ hơn nữa là khoản trừ lương không thương tiếc. May mắn thay, cô nàng Quỷ sừng sau đó đã kịp chuộc lỗi bằng cách gọi hồn đưa Robo-PON quay trở lại dương gian, dù bằng cách nào thì tôi chịu chết, không thể hiểu nổi thế giới tâm linh của họ.

Sau khi loại bỏ những món đồ bị nguyền rủa, buổi chụp ảnh diễn ra vô cùng suôn sẻ. Nghĩ lại mới thấy, quả thực cứ phải để tôi trực tiếp quản lý và điều hành thì mọi việc mới đâu vào đấy được. Trong số ảnh tôi và Haato chụp, khoảng ba phần tư đã được A-san chọn lọc cẩn thận để thiết kế và biên tập lại. Riêng tôi và Haato có giữ lại vài tấm ảnh trông ngớ ngẩn làm kỷ niệm, và dĩ nhiên đây là hàng độc quyền giới hạn, cấm lưu hành nội bộ và tuyệt đối không được bán lại đâu nhé.

Hôm nay là ngày công ty chính thức công bố bộ ảnh quảng bá. Sẵn dịp rủ hai cậu bạn Tokue và Miharu về nhà làm bài tập lớn môn lập trình, tôi đã tranh thủ giờ giải lao mở bộ ảnh lên cho cả ba cùng chiêm ngưỡng, đồng thời bàn tán rôm rả về từng thành viên xuất hiện trong đó.

\img{1-2h}

Một trong số những bức ảnh mà chúng tôi thấy buồn cười là bức mà Roboco, Matsuri, Haato và Ayame đang cưỡi trên một con gấu thật. Nàng Rô-bốt tạo một dáng vẻ rất hùng vĩ, nàng Cổ động trông hốt hoảng và bị ngã ra phía sau, Haachama thì cố bám lấy Quỷ sừng đang chực chờ rơi xuống đất.

Cầm tấm ảnh đầu tiên lên, Tokue chỉ ngón tay vào rồi quay sang hỏi tôi với vẻ mặt đầy tò mò: ``Ông bạn thấy bức này thế nào?''

``Tôi thấy nó cứ như một đống hỗn loạn ấy\dots'' tôi tặc lưỡi lắc đầu, rồi quay sang hỏi cậu bạn ngồi bên cạnh. ``Miharu, cậu thấy sao?''

Miharu nhấp một ngụm nước ngọt, thong thả đáp: ``Tớ thấy bình thường mà. Cái bức mà con bé tóc nâu này chụp riêng mới gọi là siêu phẩm hài hước ấy.''

``À, Mat-chan chứ gì?'' tôi bật cười ha hả khi nhớ lại khoảnh khắc kinh hoàng đó. ``Cái bức mà con ngươi bị trợn ngược lên như xác sống ấy! Trông vừa dị vừa buồn cười vãi chưởng!''

Tokue cũng không chịu thua, chen vào nói lớn: ``Tôi thì lại thấy bức Roboco trong trang phục con gấu đi đuổi quỷ mới là đỉnh cao nghệ thuật. Nhìn cái mặt con gấu lúc đó ngáo ngơ không chịu được!''

Nhìn đống ảnh ngộ nghĩnh này, trong đầu tôi bỗng lóe lên một ý tưởng độc đáo. Tôi đập tay xuống bàn, hào hứng đề xuất:
``Này, có khi anh em mình lấy mấy tấm ảnh này làm hình mẫu thiết kế giao diện cho bài tập code sắp tới chắc cũng ổn đấy nhỉ?''

Miharu nghe vậy liền nhìn tôi với ánh mắt đầy bất lực: ``Kuon đúng là lúc nào cũng có mấy cái ý tưởng điên rồ hết thuốc chữa.''

``Thế mà nhiều lúc đống ý tưởng điên rồ đó lại mang về điểm A cho lớp mình đấy nhé,'' Tokue cười toe toét bênh vực tôi.

Ba đứa chúng tôi cứ thế vừa ngắm nghía vừa cười đùa bình luận vui vẻ về những khoảnh khắc khó đỡ ấy. Sau đó, chúng tôi thậm chí còn táo bạo lấy chính những tấm ảnh chế này làm hình vẽ minh họa và tư liệu giao diện cho bài tập lớn nhóm của mình. Kết quả bất ngờ là sản phẩm được các thầy cô trong khoa chấm điểm cực kỳ cao! Tôi không nói đùa đâu, thầy cô còn nhận xét bài làm của nhóm rất thực tế, sáng tạo, hiệu quả và mang lại tiếng cười sảng khoái.

Xem ra, cái sự nghiệp ``làm trò hề'' mỗi ngày của các cô gái \hll\ rốt cuộc cũng mang lại chút giá trị thực tiễn giúp ích cho việc học hành của tôi đấy chứ, nhỉ? Nghĩ lại câu nói của A-san quả không sai chút nào: cái văn phòng này thực sự chẳng khác gì một rạp xiếc trung ương cả!

\end{document}